% Generated draft: requires mathematical and source review.
\documentclass[11pt]{article}
\usepackage[margin=25mm]{geometry}
\usepackage{fontspec}
\setmainfont{Latin Modern Roman}
\usepackage{amsmath,amssymb,mathtools,unicode-math}
\setmathfont{Latin Modern Math}
\usepackage{enumitem,needspace,xurl,hyperref}
\setlength{\parindent}{0pt}
\setlength{\parskip}{5pt}
\setlength{\emergencystretch}{3em}
\newcommand{\N}{\mathbb{N}}
\newcommand{\Z}{\mathbb{Z}}
\newcommand{\Q}{\mathbb{Q}}
\newcommand{\R}{\mathbb{R}}
\newcommand{\C}{\mathbb{C}}
\newcommand{\F}{\mathbb{F}}
\newcommand{\A}{\mathbb{A}}
\newcommand{\PP}{\mathbb{P}}
\newcommand{\E}{\mathbb{E}}
\newcommand{\eps}{\varepsilon}
\newcommand{\dd}{\,\mathrm{d}}
\begin{document}
\section*{UnsolvedMath ID: 1900005}
\textbf{Source problem code:} catalogue entry
\subsection{Definitions and mathematical statement}
In the usual sail model, one starts with a full-dimensional pointed cone $C\subset\R^d$ with vertex $0$, forms the convex hull of its nonzero lattice points,
\[
K(C)=\operatorname{conv}\bigl((C\cap\Z^d)\setminus\{0\}\bigr),
\]
and studies the faces of its sail, the relevant boundary of $K(C)$. The supplied record does not specify its admissible cones, any irrationality conditions, or whether an ``$m$-dimensional continued fraction'' uses ambient dimension $m$ or $m+1$. This model explains the terminology but cannot settle those conventions.

Karpenkov's Problem 5 asks which $n$-gons can occur as faces of an $m$-dimensional continued fraction. It distinguishes two regimes:
\begin{enumerate}
\item Faces at integer distance $1$ from the origin. The source describes this as the problem of describing integer convex polyhedra inscribed in full-dimensional simplices in $\R^n$.
\item Faces at integer distance greater than $1$ from the origin. The source describes this as a view-obstacle problem: lattice points that obstruct the view of the proposed face must be avoided.
\end{enumerate}
For a facet supported by a rational hyperplane $\langle a,x\rangle=b$, with primitive $a\in\Z^d$ (the coordinates have greatest common divisor $1$), its integer distance from $0$ is $|b|$. Thus the two facet regimes are $|b|=1$ and $|b|>1$. This formula does not supply an unprovided general definition for lower-dimensional faces.

A precise realizability classification must fix the permitted cones and ambient dimensions, whether $n$ denotes the number of polygon vertices or a dimension, and whether realizations are considered up to integer-affine, unimodular linear, or another equivalence. The supplied wording simultaneously uses ``$n$-gons'' and full-dimensional simplices in $\R^n$, so the repeated parameter is not silently reinterpreted. Nor does the record define ``inscribed'' or the exact view-obstacle visibility condition.

The supplied literature summary labels the problem OPEN-TRIAGE with moderate confidence. It reports partial understanding of distance-$1$ faces through empty-simplex and triangulation theory, and no resolution of the distance-$>1$ view-obstacle case in the literature it located. It supplies no exact classification theorem or supporting references for that assessment.
\subsection{Short English statement}
The question asks which lattice polygons or polyhedral faces can appear on multidimensional continued-fraction sails, separating integer distance $1$ from greater distances. The source leaves dimension conventions, equivalence of realizations, and visibility constraints unspecified. Its summary reports only partial understanding of the first regime.
\subsection{Sources}
\begin{itemize}
\item Local UnsolvedMath registry, permanent ID 1900005, exact source code catalogue entry, ``Geometry of Continued Fractions --- Faces of sails''; status open, classification OPEN-TRIAGE.
\item Karpenkov, \emph{Open problems in geometry of continued fractions} (2017), Problem 5, \url{https://arxiv.org/abs/1712.01450}. The record identifies extraction from public arXiv source TeX, accessed July 29, 2026, and says the source presented the problem as open.
\item The supplied research summary, with moderate confidence, mentions empty-simplex/triangulation theory for distance-$1$ faces and an unresolved distance-$>1$ view-obstacle case. No independent authoritative citation or precise partial theorem is supplied for those claims.
\end{itemize}
\end{document}
